% Chapter introduction: why the topic arose, who first asked, what it answers.
% STATUS: written 2026-11-05 (cites only the chapter's own sources).
% Read in full: Singer-Wermer 1955, Kadison 1966, Sakai 1960, Sakai 1966,
% Sakai 1976, Johnson-Sinclair 1968, Ringrose 1972, Han 1993, Lin 1971,
% Williams 1971, Dash 1975, Hyers 1941, Miura-Hirasawa-Takahasi 2006;
% Sakai 1991 preliminaries and contents.
% Additions of 2026-11-07 name at the cited pages: Bade-Curtis 1974,
% Mathieu 1994, Molnar 1994, Niknam 1983, Semrl 1991, Ali 2011,
% Kadison-Liu 2014, Yu-Zhang 2012, Liu-Hou 2025, Bouhafsi-Ech-chad-Zouaki
% 2022, Roh-Jung-Chang 2009, Konig-Milman 2014, Hosseini-Peralta-Su 2025,
% Eisaei-Moghimi 2024.
\v{S}ilov proved in 1947 that the algebra of all infinitely differentiable
functions on an interval cannot be normed so as to be a Banach algebra
\cite[p.~260]{SingerWermer1955}. Kaplansky proposed that the reason was
a theorem about derivations: a commutative semisimple Banach algebra should
carry no nonzero derivation \cite[p.~260]{SingerWermer1955}. Singer and
Wermer proved the conjecture for bounded derivations in 1955
\cite[Theorem~1, p.~260]{SingerWermer1955}. Their paper opened the study of
derivations on Banach and operator algebras. Kaplansky then asked whether a
derivation of a $C^*$-algebra is automatically continuous. Sakai settled that
question in 1960 \cite[p.~31]{Sakai1960Kaplansky}. The same circle of
ideas led Kadison and Sakai in 1966 to the theorem that every derivation
of a von Neumann algebra is inner \cite[Theorem~1, p.~274]{Sakai1966}.

The chapter tells this story in the order of the questions.
Section~\ref{sec:banach-operator:singer-wermer} reports the commutative
theory of Singer and Wermer and its completion by Thomas.
Section~\ref{sec:banach-operator:kadison-sakai} gives the Kadison--Sakai
theorem on von Neumann algebras and the spatial theorem of Kadison for
$C^*$-algebras on a Hilbert space. It adds the extended derivations of the
Murray--von Neumann algebra of a finite von Neumann algebra
\cite[Theorem~5, p.~2090]{KadisonLiu2014Murray} and the nonlinear Lie
derivations of a factor \cite[Theorem~2.1, p.~1980]{YuZhang2012StarLie}. It
closes with the 2-local derivations
of Ayupov, Arzikulov and Kudaybergenov, which turn out to be derivations on
every von Neumann algebra \cite[Theorem~5, p.~7]{AyupovArzikulov2011TwoLocal}.
It also gives the local derivations of Hadwin, Li, Li and Ma on measurable
operators \cite[Theorem~1, p.~2]{HadwinLiLiMa2013}. On the measurable
operators themselves the derivations that are linear over the centre are
spatial \cite[Theorem~3.9, p.~11]{AlbeverioAyupovKudaybergenov2007}.
Section~\ref{sec:banach-operator:automatic-continuity} explains why
derivations of semisimple Banach algebras are continuous, the theorem of
Johnson and Sinclair, and the module form of Ringrose. It adds the
continuity ideals of Bade and Curtis
\cite[Corollary~2.8, p.~379]{BadeCurtis1974}. It adds too the continuity of
strongly generalized derivations
\cite[Theorem~2.7, p.~5]{HosseiniPark2023}, and the failure of inner
derivations to extend from a dense ideal
\cite[Theorem~4.1, p.~3]{ShafieaslTavakkoli2026}.
Section~\ref{sec:banach-operator:unbounded} passes to unbounded
derivations, the generators of one-parameter automorphism groups in quantum
dynamics, as Sakai began the theory. It records the kernel
stabilization of weakly defined derivations
\cite[Theorem~3.5, p.~7]{Ismert2018Kernel}. It records too that a derivation
with a range that is not dense can be closed
\cite[Theorem~1, p.~181]{Niknam1983Closable}.
Section~\ref{sec:banach-operator:operator-derivations} moves to nest
algebras and additive derivations, the work of Christensen and Han. It adds
the Lie-type derivations of Zhang and Wei
\cite[Theorem~3.1, p.~6]{ZhangWei2017}, the additive derivations of a full
operator algebra \cite[Theorem~2.3, p.~236]{Semrl1991Additive} and the Lie
derivations at a fixed product \cite[Theorem~2.2, p.~81]{LiuHou2025Lie}.
Section~\ref{sec:banach-operator:commutator-ranges} studies the range of a
commutator map, from Kleinecke and Shirokov to the
Wintner--Wielandt--Halmos statement. It adds the ranges of normal Jordan
$*$-derivations \cite[Theorem~1, p.~692]{Molnar1994Range} and the
$P$-symmetric pairs of Bouhafsi, Ech-chad and Zouaki
\cite[Theorem~3.4, p.~149]{BouhafsiEchChadZouaki2022}.
Section~\ref{sec:banach-operator:stability} closes with approximate
derivations and the Hyers--Ulam question of whether an almost derivation is
near an exact one. It carries the question to Jordan triple higher
derivations \cite[Theorem~2.6, p.~1984]{RohJungChang2009Triple} and to the
derivation equation for $C^k$-functions
\cite[Theorem~1, p.~407]{KonigMilman2014Derivation}.

\hbfigure{18-banach-operator/unbounded}{On $C[0,2\pi]$ with the sup norm, $\sin nx$ has norm $1$ and its derivative has norm $n$. Differentiation is an unbounded derivation, here $n=4$.}{fig:banach:unbounded}
